\thispagestyle{empty}
\begin{center}
    \vspace*{1.5cm}
    
    {\fontsize{28}{34}\bfseries\color{agriGreen} จุลชีววิทยาสำหรับการเกษตร}\\[10pt]
    {\fontsize{18}{22}\bfseries\color{agriLeaf} (Microbiology for Agriculture)}\\[15pt]
    {\color{agriGold}\rule{0.6\textwidth}{2pt}}\\[20pt]
    
    {\fontsize{14}{18}\color{agriSlate} ตำราประกอบการเรียนการสอนและการวิจัย}\\[8pt]
    {\fontsize{12}{16}\color{slateGrey} ตามกรอบมาตรฐานคุณวุฒิระดับอุดมศึกษาแห่งชาติ (TQF)}\\[60pt]
    
    {\fontsize{14}{18}\bfseries\color{agriGreen} ผู้เขียน}\\[8pt]
    {\fontsize{20}{26}\bfseries\color{agriEarth} ผู้ช่วยศาสตราจารย์ ดร.จิรภัทร จันทมาลี}\\[10pt]
    {\fontsize{12.5}{16}\color{agriSlate} วท.บ. (จุลชีววิทยา เกียรตินิยมอันดับ 2), วท.ม. (จุลชีววิทยาอุตสาหกรรม), Ph.D. (Microbiology) จุฬาฯ}\\[6pt]
    {\fontsize{13}{16}\color{agriSlate} อาจารย์ประจำสาขาวิชาจุลชีววิทยา}\\[4pt]
    {\fontsize{13}{16}\color{agriSlate} คณะวิทยาศาสตร์และเทคโนโลยี มหาวิทยาลัยราชภัฏรำไพพรรณี}\\[80pt]
    
    \vfill
    
    {\fontsize{13}{16}\bfseries\color{agriGreen} คณะวิทยาศาสตร์และเทคโนโลยี มหาวิทยาลัยราชภัฏรำไพพรรณี}\\[6pt]
    {\fontsize{12}{15}\color{agriSlate} พุทธศักราช 2569}
\end{center}
\clearpage
